\documentclass[a4paper,twoside]{article}

\usepackage[T1]{fontenc}
\usepackage{epsfig}
\usepackage{subcaption}
\usepackage{calc}
\usepackage{amssymb}
\usepackage{amstext}
\usepackage{amsmath}
\usepackage{amsthm}
\usepackage{amsfonts}
\usepackage{array}
\usepackage{booktabs}
\usepackage{color}
\usepackage{multicol}
\usepackage{pslatex}
\usepackage{apalike}
\usepackage[ruled,vlined]{algorithm2e}
\usepackage[bottom]{footmisc}
\usepackage[hidelinks]{hyperref}
\usepackage{SCITEPRESS}

\newtheorem{definition}{Definition}
\newtheorem{proposition}{Proposition}
\newtheorem{remark}{Remark}

\renewcommand{\T}{\mathcal{T}}
\renewcommand{\R}{\mathbb{R}}
\newcolumntype{L}[1]{>{\raggedright\arraybackslash}p{#1}}

\sloppy

\title{Preference-Aware Overlapping De-Reservation under Score Constraints}

% Double-blind review version. Replace with real author data for camera-ready submission.
\author{\authorname{Anonymous Author(s)}
\affiliation{\sup{1}Anonymous Institution}
\email{anonymous@example.com}
}

\begin{document}

\abstract{We study admissions with one open category and multiple reserve categories when students may qualify for several reserve categories but, under one-to-one counting, can be assigned to only one. We impose a score constraint requiring each active reserve cutoff to stay within a fixed gap of the open cutoff. We present a compact formulation of \emph{Preference-Aware Overlapping De-Reservation} (PA-ODR), related to de-reservation methods in reserve-based admissions: it fills the open category by raw scores, matches reserves only through score-feasible student-category pairs, and de-reserves seats that cannot be filled feasibly. Within the feasible region, category preferences are used only after maximizing match size and score performance. We also record a simple impossibility: score-feasibility, non-wastefulness, and a stronger open-merit requirement may fail to coexist even in a minimal instance.}

\keywords{Matching, Affirmative Action, Reserve Categories, De-reservation, School Choice}

\onecolumn \maketitle \normalsize \setcounter{footnote}{0} \vfill

\section{\uppercase{Introduction}}

Reserve systems are often described category by category, but many admissions environments are not category-separable. A student may qualify for more than one protected or priority group, while the institution may still require one-to-one counting: each admitted student fills at most one seat category. This makes the assignment of an overlapping student to a particular reserve category a substantive part of the mechanism rather than a bookkeeping detail.

The issue becomes sharper when admissions are also subject to a cutoff rule. In the Indian affirmative-action debate, for example, cutoff gaps between open and reserve admissions have been treated as an institutional constraint rather than merely as an outcome statistic. A reserve seat may then be legally or politically available but unusable at the realized open cutoff. In such a case, ordinary non-wastefulness is too coarse: an empty reserve seat is not necessarily wasted if every available eligible student would violate the score-gap rule.

This short paper focuses on one constructive response. Preference-Aware Overlapping De-Reservation (PA-ODR) treats the score constraint as hard. It first determines the open cutoff from raw scores, then assigns reserve seats only through student-category pairs that satisfy the cutoff rule. If a reserve category cannot be filled under the current cutoff, the mechanism de-reserves one unit of its capacity and transfers it to the open category. Overlap is handled by a matching problem on the feasible reserve graph, with students' category preferences used only after match size and score performance have been fixed.

The contribution is deliberately narrower than a full characterization. We give a compact model, the mechanism, and the basic guarantees: termination, score-feasibility, non-wastefulness relative to the final capacity vector, and within-category fairness. We also record why a reserve mechanism that insists on fixed capacities and open merit can run into a feasibility conflict. The point is not that de-reservation is costless; rather, it makes explicit which reserve capacity cannot be used while respecting the score constraint.

\section{\uppercase{Model}}

Let $I$ be a finite set of students, $T$ a finite set of reserve categories, $o$ the open category, and $V=T\cup\{o\}$. Each category $c\in V$ has capacity $q^c\in\mathbb{N}$. Each student $i\in I$ has raw score $\sigma(i)\in\R_+$, reserve-eligibility set $\T(i)\subseteq T$, and a strict preference ordering $\succ_i$ over $V_i:=\{o\}\cup\T(i)$. This combines overlapping types familiar from controlled school choice and multidimensional affirmative action \cite{KHIY17a,AyBo21a} with a cutoff-based reserve restriction.

An assignment is a tuple $\mu=(\mu_c)_{c\in V}$, where $\mu_c\subseteq I$ is the set of students assigned to category $c$.

\begin{definition}[Feasibility]
\label{def:paodr-feasibility}
An assignment $\mu$ is feasible if $|\mu_c|\le q^c$ for every category $c\in V$ and $\mu_c\cap\mu_{c'}=\emptyset$ for every distinct pair $c,c'\in V$.
\end{definition}

The second condition is one-to-one counting. For each reserve category $t\in T$, let $k_t\ge 0$ be its score bonus. The open category uses raw scores, while reserve category $t$ evaluates student $i$ by the adjusted score $\sigma_t(i)=\sigma(i)+k_t$ whenever $t\in\T(i)$.

For a feasible assignment $\mu$, define the open cutoff by $e_o(\mu):=\min_{i\in\mu_o}\sigma(i)$ when $\mu_o\neq\emptyset$, and the adjusted reserve cutoff by $e_t^+(\mu):=\min_{i\in\mu_t}\sigma_t(i)$ when $\mu_t\neq\emptyset$. Let $\delta\ge 0$ be the maximum allowed score gap.

\begin{definition}[Score-feasibility]
\label{def:paodr-score-feasibility}
A feasible assignment $\mu$ is score-feasible if for every reserve category $t\in T$ with $\mu_o\neq\emptyset$ and $\mu_t\neq\emptyset$,
\[
e_o(\mu)-e_t^+(\mu)\le \delta.
\]
\end{definition}

This is an outcome restriction on final cutoffs: once reserve bonuses are treated as institutionally relevant, reserve admissions should not fall too far below open admissions. It is therefore separate from the ranking rule used inside reserve categories.

\subsection{Desiderata}

We use three basic desiderata. First, a feasible assignment is non-wasteful relative to a capacity vector if no unmatched student can be added to an eligible category without violating capacities and score-feasibility. For PA-ODR the relevant vector is the final effective capacity vector $\hat q$, because de-reservation changes which seats remain reserve seats.

Second, within-category fairness requires that no unmatched eligible student with a higher adjusted score be blocked from a reserve category by a lower adjusted-score assignee. This condition is intentionally local to reserve categories. It is weaker than full no justified envy, because a student assigned through one route may prefer another route, and PA-ODR does not promise to satisfy such route preferences globally.

Third, score-feasibility is treated as a hard institutional constraint. This is the main difference between PA-ODR and a direct deferred-acceptance construction such as SPSE. The latter lets students propose to categories and lets categories hold their highest adjusted-score applicants, but the resulting reserve cutoffs may still fall too far below the open cutoff.

\section{\uppercase{PA-ODR}}

The mechanism treats score-feasibility as a hard constraint. It preserves total capacity while allowing infeasible reserve seats to be transferred to the open category.

For a fixed open cutoff $e$, define a bipartite graph $G(e)$ between students not assigned to the open category and reserve categories. An edge $(i,t)$ is present exactly when
\[
t\in\T(i)\quad\text{and}\quad \sigma(i)+k_t\ge e-\delta.
\]
Thus a student may be feasible for one reserve category but not for another.

Among all capacitated matchings on $G(e)$, PA-ODR uses the following lexicographic rule: first maximize the number of matched reserve pairs; then lexicographically maximize matched adjusted scores; finally lexicographically minimize students' category-preference ranks.

This reserve-side step can be implemented as a maximum-weight bipartite matching with sufficiently separated weights. The first component rewards cardinality, the second rewards adjusted scores, and the third penalizes assignments to less-preferred categories. The exact numerical encoding is not conceptually important here; what matters is that overlap is solved jointly across categories rather than by processing categories independently. Independent processing can assign a multiply eligible student to one category while leaving another category infeasible, even when a different routing would satisfy the score rule.

\begin{algorithm*}[t]
\caption{Preference-Aware Overlapping De-Reservation (PA-ODR)}
\label{alg:paodr-short}
\KwIn{Students $I$, scores $\sigma$, eligibility $\T(\cdot)$, preferences $(\succ_i)$, capacities $q^o,(q^t)_{t\in T}$, bonuses $(k_t)_{t\in T}$, gap $\delta$}
\KwOut{Assignment $\mu=(\mu_o,(\mu_t)_{t\in T})$}
Initialize effective capacities $\hat q^o\gets q^o$ and $\hat q^t\gets q^t$ for all $t\in T$\;
\While{true}{
    Assign the top $\min\{\hat q^o,|I|\}$ raw-score students to $o$; call this set $\mu_o$\;
    Let $e=\min_{i\in\mu_o}\sigma(i)$ if $\mu_o\neq\emptyset$, and $e=-\infty$ otherwise\;
    Build $G(e)$ with edge $(i,t)$ iff $t\in\T(i)$ and $\sigma(i)+k_t\ge e-\delta$\;
    Compute the reserve assignment $(\mu_t)_{t\in T}$ by the lexicographic rule, with capacities $\hat q^t$\;
    \eIf{there exists $t$ with $|\mu_t|<\hat q^t$}{
        Choose one such category $t^\ast$\;
        Update $\hat q^{t^\ast}\gets \hat q^{t^\ast}-1$ and $\hat q^o\gets \hat q^o+1$\;
    }{
        \Return{$\mu=(\mu_o,(\mu_t)_{t\in T})$}\;
    }
}
\end{algorithm*}

The mechanism preserves the open-first merit order and uses de-reservation only when a reserve seat cannot be filled without violating the score constraint.

\begin{remark}[What de-reservation means here]
De-reservation is not a claim that the original reserve category was normatively unimportant. It is an operational response to a binding cutoff constraint. If the current open cutoff makes a reserve seat unfillable by any score-feasible eligible student, PA-ODR converts that unit into open capacity rather than leaving the mechanism with an infeasible reserve assignment.
\end{remark}

\section{\uppercase{Properties}}

\begin{proposition}[Termination]
\label{prop:paodr-termination-short}
PA-ODR terminates after finitely many iterations.
\end{proposition}

\begin{proof}
Whenever the mechanism does not stop, it reduces the effective capacity of some reserve category by one. Hence the total effective reserve capacity $\sum_{t\in T}\hat q^t$ strictly decreases at every non-terminating iteration and cannot decrease forever.
\end{proof}

\begin{proposition}[Score-feasibility]
\label{prop:paodr-score-short}
The output of PA-ODR is score-feasible.
\end{proposition}

\begin{proof}
Fix the final iteration and let $e$ be its open cutoff. Every reserve assignee $i\in\mu_t$ was selected along an edge of $G(e)$, so $\sigma(i)+k_t\ge e-\delta$. Therefore $e_t^+(\mu)\ge e-\delta=e_o(\mu)-\delta$, which implies $e_o(\mu)-e_t^+(\mu)\le \delta$ for each nonempty reserve category.
\end{proof}

\begin{proposition}[Non-wastefulness with respect to the final capacity vector]
\label{prop:paodr-nonwaste-short}
The output of PA-ODR is non-wasteful relative to the final effective capacity vector $\hat q$.
\end{proposition}

\begin{proof}
At termination, every reserve category with positive effective capacity is full; otherwise the mechanism would de-reserve another unit and continue. The open category is also filled up to $\hat q^o$ unless there are fewer students than total capacity. Hence no unmatched student can be added while preserving the final capacities and score-feasibility.
\end{proof}

\begin{proposition}[Within-category fairness]
\label{prop:paodr-within-short}
The output of PA-ODR is fair within categories.
\end{proposition}

\begin{proof}
Fix the final iteration and its open cutoff $e$. Suppose an unmatched student $i$ could feasibly replace some $j\in\mu_t$ in a reserve category $t$ with a strictly higher adjusted score. Then $(i,t)$ is an edge of the final graph $G(e)$, and replacing $j$ with $i$ keeps the reserve match size unchanged while lexicographically improving the vector of matched adjusted scores. This contradicts the score-maximality of the chosen reserve matching.
\end{proof}

\section{\uppercase{SPSE in Brief}}

The student-proposing score-elevated mechanism (SPSE) is the deferred-acceptance benchmark in the longer draft. Students apply to eligible categories in preference order; each category tentatively keeps its highest adjusted-score applicants up to capacity. SPSE keeps the initial capacity vector $q$ fixed and does not directly enforce the score-gap rule in Definition~\ref{def:paodr-score-feasibility}.

SPSE is attractive because it has the familiar logic of student-proposing deferred acceptance. It respects one-to-one counting because a student can be tentatively held by at most one category, and it gives students a clear proposal order over admissible categories. The weakness is that score-feasibility is not built into the rejection rule. A reserve category may hold its best applicants and still end with a cutoff below the allowed gap from the open cutoff.

This distinction is useful for interpreting PA-ODR. PA-ODR is not meant to dominate SPSE on every criterion. It gives up the fixed capacity vector and can create justified envy over routes. Its role is different: when the cutoff rule is non-negotiable, PA-ODR supplies a feasible assignment by changing the effective capacities rather than by weakening the score constraint.

\begin{proposition}[A simple impossibility under stronger fairness]
\label{prop:simple-impossibility-short}
There need not exist any matching that simultaneously satisfies feasibility, non-wastefulness, score-feasibility, within-category fairness, and open-merit, where open-merit means that no unmatched student has a higher raw score than an open-category assignee.
\end{proposition}

\begin{proof}
Consider one open category $o$ and one reserve category $t$, each with capacity $1$. Let $\delta=10$ and $k_t=0$. Let student $A$ have score $100$ and no reserve eligibility, and let student $B$ have score $80$ and eligibility $\{t\}$. The matching $(\{A\},\{B\})$ violates score-feasibility because $100-80=20>10$. The matching $(\{B\},\emptyset)$ violates open-merit because $A$ is unmatched and has a higher raw score than $B$. The matching $(\{A\},\emptyset)$ violates non-wastefulness because $B$ can be added feasibly to $t$. The matching $(\emptyset,\{B\})$ also violates non-wastefulness because $A$ can be added feasibly to $o$. The empty matching violates non-wastefulness as well. Hence no feasible matching satisfies all five requirements simultaneously.
\end{proof}

\section{\uppercase{Comparison with SPSE}}

SPSE keeps the initial reserve capacities and delivers deferred-acceptance-style fairness properties, but it need not enforce the score-gap rule. PA-ODR instead enforces score-feasibility by construction while allowing the effective reserve capacities to change through de-reservation.

\begin{center}
\textbf{Table 1:} Comparison of SPSE and PA-ODR.

\small
\renewcommand{\arraystretch}{1.15}
\begin{tabular}{|L{0.18\textwidth}|L{0.17\textwidth}|L{0.20\textwidth}|}
\hline
\textbf{Property} & \textbf{SPSE} & \textbf{PA-ODR} \\
\hline
One-to-one counting & Yes & Yes \\
Feasibility & Yes & Yes \\
Non-wastefulness & Yes & Yes, w.r.t.\ final $\hat q$ \\
Score-feasibility & Not guaranteed & Guaranteed \\
Capacity reallocation & Fixed at $q$ & Endogenous \\
Preferences over categories & Via proposals & Final tie-breaker \\
\hline
\end{tabular}
\end{center}

PA-ODR therefore isolates a central modeling choice: enforcing the score-gap rule can naturally be combined with endogenous capacity reallocation, especially once the stronger open-merit requirement of Proposition~\ref{prop:simple-impossibility-short} is taken seriously.

\section{\uppercase{Discussion}}

The main limitation of PA-ODR is that the amount of de-reservation can depend on the order in which underfilled categories are processed. This order dependence should be specified transparently before the mechanism is used. A natural alternative is to search for a minimum-de-reservation capacity vector, but that changes the problem from a simple iterative rule into an optimization problem over capacity vectors.

A second limitation is incentives. Because PA-ODR uses students' category preferences in the reserve-side tie-breaker, and because de-reservation changes the open cutoff endogenously, it is not a strategy-proof mechanism in the same sense as student-proposing deferred acceptance. This is the price of making score-feasibility binding while still using preferences to route overlapping students.

The mechanism is therefore best read as a policy fallback for environments where the cutoff rule is hard and the institution is willing to let infeasible reserve capacity move to the open category. In that setting, PA-ODR gives a transparent allocation procedure and separates three questions that are often conflated: who is eligible for which reserve categories, which assignments satisfy the score rule, and which reserve capacity can actually be used at the final open cutoff.

\bibliographystyle{apalike}
\bibliography{references}

\end{document}
